\chapter{Теоретические сведения}
\label{ch:theory}

\section{Система доменных имён DNS}

\textbf{DNS} (Domain Name System) --- иерархическая распределённая система именования
узлов сети. Пространство имён имеет древовидную структуру: корень, домены верхнего
уровня (\texttt{.ru}, \texttt{.com}, \texttt{.local}), домены второго уровня и т.д.

Основные понятия:
\begin{itemize}
  \item \textbf{Зона} --- административная единица пространства имён;
  \item \textbf{Авторитативный сервер} --- хранит оригинальные записи зоны;
  \item \textbf{Рекурсивный резолвер} --- выполняет цепочку запросов от имени клиента;
  \item \textbf{Форвардер} --- проксирует запросы для нелокальных зон на внешний сервер.
\end{itemize}

\section{Типы ресурсных записей DNS}

Основные типы записей, используемых в работе:
\begin{itemize}
  \item \textbf{SOA} (Start of Authority) --- первичный NS, e-mail администратора,
        параметры обновления зоны;
  \item \textbf{NS} (Name Server) --- авторитативный сервер зоны;
  \item \textbf{A} --- адресная запись: имя хоста $\to$ IPv4;
  \item \textbf{PTR} (Pointer) --- обратная запись: IP-адрес $\to$ имя хоста;
  \item \textbf{CNAME} (Canonical Name) --- псевдоним;
  \item \textbf{MX} (Mail Exchanger) --- почтовый сервер домена.
\end{itemize}

\section{Прямая и обратная зоны}

\textbf{Прямая зона} (\textit{forward zone}) содержит записи типов SOA, NS, A.
Имя зоны совпадает с доменом, например \texttt{STUDENT.GROUP.local}.

\textbf{Обратная зона} (\textit{reverse zone}) строится инвертированием октетов
подсети с суффиксом \texttt{.in-addr.arpa}. Для сети
\texttt{192.168.\cfgLabStudentN.0/24} имя зоны:
\texttt{\cfgLabStudentN.168.192.in-addr.arpa}. Содержит записи SOA, NS, PTR.

\section{BIND9}

\textbf{BIND9} (Berkeley Internet Name Domain) --- наиболее распространённое ПО
DNS-сервера. Устанавливается командой:
\begin{lstlisting}
apt install bind9 dnsutils
\end{lstlisting}

Конфигурационные файлы в \texttt{/etc/bind/}:
\begin{itemize}
  \item \texttt{named.conf} --- главный файл, подключает остальные;
  \item \texttt{named.conf.options} --- форвардеры, \texttt{listen-on}, dnssec;
  \item \texttt{named.conf.local} --- описание локальных зон;
  \item файлы зон --- в \texttt{/var/lib/bind/} (BIND имеет права на запись,
        необходимо для DDNS).
\end{itemize}

Пример блока \texttt{named.conf.options}:
\begin{lstlisting}
options {
    directory "/var/cache/bind";
    forwarders { 8.8.8.8; };
    dnssec-validation auto;
    auth-nxdomain no;
    listen-on { 127.0.0.1; 192.168.N.1; };
};
\end{lstlisting}

\texttt{forwarders} --- DNS Google для внешних запросов; \texttt{listen-on} ---
принимать запросы только с localhost и LAN-интерфейса.

\section{Динамическое обновление DNS (DDNS)}

\textbf{DDNS} (RFC~2136) позволяет DHCP-серверу автоматически создавать
A- и PTR-записи при выдаче IP-адреса клиенту. Для аутентификации используется
TSIG-ключ \texttt{rndc-key} (HMAC-MD5), хранящийся в \texttt{/etc/bind/rndc.key}.

Этапы настройки DDNS:
\begin{enumerate}
  \item Скопировать ключ RNDC в каталог DHCP:
        \texttt{cp -Rr /etc/bind/rndc.key /etc/dhcp/ddns-keys/}
  \item Добавить в \texttt{/etc/dhcp/dhcpd.conf} директивы
        \texttt{ddns-updates on}, \texttt{ddns-update-style standard},
        описания зон и ссылку на ключ.
  \item Перезапустить bind9 и isc-dhcp-server.
  \item Проверить регистрацию клиента через \texttt{nslookup}.
\end{enumerate}

После подключения Desktop-клиента DHCP выдаёт ему IP и одновременно
регистрирует запись \texttt{HOSTNAME.STUDENT.GROUP.local} в DNS.
Журнальные файлы \texttt{*.jnl} в \texttt{/var/lib/bind/} хранят историю
динамических обновлений.

\endinput
